\chapter{\textit{Moringa} compound evidence and candidate prioritisation}\label{ch:compounds}

\section{Verified constituents of \textit{Moringa oleifera} seed and leaf}\label{sec:constituents}

Curation produced a provenance-carrying inventory of \textit{M.~oleifera} seed and leaf
constituents drawn from primary analytical studies, with the extraction solvent,
identification method, DOI verification route and reading depth of each report recorded
(Table~\ref{tab:c1}). Three findings are substantive in themselves.

\paragraph{Only one constituent is confirmed in the solvent the original study used.}
\textcite{abdelsayed2021jfda} isolated \textbf{niazimicin}, a thiocarbamate glycoside, from
the \textbf{seed methylene-chloride fraction}, confirmed its structure by HRMS, UV, IR, EI-MS
and one- and two-dimensional NMR, and quantified it by validated UPLC-MS/MS in both the seed
and the leaf \ce{CH2Cl2} fractions. The same study did \emph{not} detect niazirin or
glucomoringin there. The only other dichloromethane-specific identifications are GC-MS
spectral-library matches in leaf DCM material, read at abstract level:
4-hydroxyphenylacetonitrile and methyl 4-hydroxyphenylacetate \parencite{mohdfisall2021},
and palmitic acid and \textit{cis}-vaccenic acid \parencite{khandelwal2019}. The latter
appeared in a journal indexed in neither PubMed nor Europe PMC, its full compound list could
not be accessed, and electron-impact library matching cannot distinguish the $\Delta^{9}$
and $\Delta^{11}$ C18:1 isomers without authentic standards.

\paragraph{The identification of several widely cited \enquote{Moringa compounds} is weaker
than their reputation.} Free quercetin and kaempferol rest only on tentative LC-HRMS
dereplication without authentic standards and may form by hydrolysis of the abundant
glycosides during extraction; the forms that primary analytical work actually establishes in
leaf are the \textbf{glycosides} isoquercitrin, astragalin and rutin. Under the prespecified
rules the tentative aglycone reports still count as primary seed or leaf reports, so
quercetin and kaempferol \emph{are} assessable, but only in the weakest identification band
($D_1=1$). The pharmacokinetic caution applies in both directions: aglycone evidence may not
be transferred to a glycoside, or the reverse, without explicit justification. Benzyl
isothiocyanate is \textbf{not} a documented seed or leaf constituent: its precursor, benzyl
glucosinolate, is reported in \emph{roots}, while seed and leaf carry the
rhamnosyloxy-benzyl glucosinolate instead. Gallic acid and caffeic acid yielded no primary
report in \textit{M.~oleifera} seed or leaf and do not appear in the inventory.

\paragraph{Identifier hygiene was not a formality.} Four public records returned by a name
search are wrong or ambiguous for this purpose, and using them would have propagated an
error: a mis-structured carbamate record (\ce{C16H23NO8}, with the ethoxy group on the
sugar); a \enquote{niaziminin A} record whose structure is not the thiocarbamate; a
niazimicin record without stereochemistry; and glucomoringin, for which every available
record is a potassium salt with a different stereo-InChIKey and none could be matched to the
natural structure. The glucomoringin identifier remains \textbf{unresolved}.

\section{The candidate shortlist}\label{sec:shortlist}

The evidence linking these constituents to radiation-response biology was curated
separately, from primary studies of purified compounds only (Table~\ref{tab:c2}), and the
prespecified score of Equation~\ref{eq:priority} was then applied
(Table~\ref{tab:ranking}; Figures~\ref{fig:priority} and~\ref{fig:evidence-map}).
\textbf{Eighteen compounds entered the ranked shortlist and five reached tier~A}; three
further constituents scored below the tier-C threshold and were not ranked, and benzyl
isothiocyanate was not assessable. The written justification of every point is given in
Table~\ref{tab:c3reasons}.

\begin{table}[tbp]
  \caption[Candidate ranking]{Candidate ranking of documented \textit{M.~oleifera} seed and
    leaf constituents under the prespecified four-dimension score (Equation~\ref{eq:priority};
    $D_1$ identity 0--2, $D_2$ radiation biology 0--3, $D_3$ independent sources 0--2, $D_4$
    feasibility 0--1). Extraction evidence is shown beside the score and is not part of it.
    Rows follow the ranking file. Every tier is exploratory and hypothesis-generating.
    Source: \path{results/tables/primary_v1/compounds_v1/table_c3_candidate_ranking.csv}.}
  \label{tab:ranking}
  \footnotesize\setlength{\tabcolsep}{3.4pt}
  \begin{tabular}{@{}r L{0.235\linewidth} l l c c c c c L{0.10\linewidth} L{0.15\linewidth}@{}}
    \toprule
    \thead{Rank} & \thead{Compound} & \thead{Part} & \thead{Extraction}
      & \thead{$D_1$} & \thead{$D_2$} & \thead{$D_3$} & \thead{$D_4$} & \thead{Total}
      & \thead{Tier} & \thead{Liability} \\
    \midrule
    1 & $\alpha$-Tocopherol & both & other solvent & 2 & 3 & 2 & 1 & 8 & \textbf{A} & \na \\
    2 & Chlorogenic acid & leaf & other solvent & 1 & 3 & 2 & 1 & 7 & \textbf{C} (capped from A)
      & protects tumour cells (CGA-01) \\
    3 & Kaempferol & both & indirect & 1 & 3 & 2 & 1 & 7 & \textbf{A} & \na \\
    4 & Quercetin & both & indirect & 1 & 3 & 2 & 1 & 7 & \textbf{A} & \na \\
    5 & Rutin & leaf & other solvent & 1 & 3 & 2 & 1 & 7 & \textbf{A} & \na \\
    6 & $\beta$-Sitosterol & seed & other solvent & 1 & 3 & 2 & 1 & 7 & \textbf{A} & \na \\
    \na & Benzyl isothiocyanate & other & indirect & 0 & 3 & 2 & 1 & 6 & not assessable & \na \\
    7 & Niazirin & both & other solvent & 2 & 1 & 2 & 1 & 6 & B & \na \\
    8 & Oleic acid & seed & other solvent & 1 & 2 & 1 & 1 & 5 & B & \na \\
    9 & Astragalin & leaf & other solvent & 1 & 1 & 2 & 1 & 5 & C (capped from B)
      & genotoxic (AST-01) \\
    10 & Moringin & both & other solvent & 1 & 1 & 2 & 1 & 5 & B & \na \\
    11 & Phytol & leaf & other solvent & 1 & 1 & 2 & 1 & 5 & B & \na \\
    12 & Glucomoringin & both & other solvent & 2 & 0 & 2 & 0 & 4 & C & \na \\
    13 & Niazimicin & both & \textbf{DCM} & 2 & 0 & 2 & 0 & 4 & C & \na \\
    14 & Acetylated moringin\textsuperscript{a} & leaf & other solvent & 2 & 0 & 1 & 0 & 3 & C & \na \\
    15 & Niaziminin & leaf & other solvent & 2 & 0 & 1 & 0 & 3 & C & \na \\
    16 & \textit{O}-Ethyl rhamnosyloxybenzyl carbamate\textsuperscript{b} & seed & other solvent
      & 2 & 0 & 1 & 0 & 3 & C & \na \\
    17 & Isoquercitrin & leaf & other solvent & 1 & 0 & 1 & 0 & 2 & C & \na \\
    18 & Palmitic acid & both & \textbf{DCM} & 1 & 0 & 1 & 0 & 2 & C & \na \\
    \na & 4-Hydroxyphenylacetonitrile & both & \textbf{DCM} & 1 & 0 & 0 & 0 & 1 & not prioritised & \na \\
    \na & Methyl 4-hydroxyphenylacetate & leaf & \textbf{DCM} & 1 & 0 & 0 & 0 & 1 & not prioritised & \na \\
    \na & \textit{cis}-Vaccenic acid & leaf & \textbf{DCM} & 1 & 0 & 0 & 0 & 1 & not prioritised & \na \\
    \bottomrule
  \end{tabular}
  \tablenote{Extraction: \textbf{DCM}, identified in a dichloromethane extract or fraction of
    seed or leaf; other solvent, identified in seed or leaf with another solvent; indirect,
    identification indirect or unconfirmed. For benzyl isothiocyanate the plant part is
    \enquote{other} because its only \textit{Moringa} report concerns its precursor in roots.
    \textsuperscript{a}\,4-[(4$'$-\textit{O}-Acetyl-$\alpha$-L-rhamnosyloxy)benzyl]
    isothiocyanate. \textsuperscript{b}\,\textit{O}-Ethyl-4-($\alpha$-L-rhamnosyloxy)benzyl
    carbamate.}
\end{table}

Three features of the ranking matter more than the order itself.

\paragraph{The liability cap did real work.} Chlorogenic acid scored 7 of 8, equal to four
tier-A compounds, and is nevertheless held at tier~C because one primary study reports that
it \emph{protects hepatocellular carcinoma cells} from radiation. A score alone would have
promoted it; the prespecified cap prevented that.

\paragraph{The exclusion rule did real work.} Benzyl isothiocyanate scored 6 on biology and
sources and is still \emph{not assessable}, because it has no primary report as a seed or leaf
constituent. It cannot be ranked however strong its radiation evidence is.

\paragraph{The mismatch exposed by the curation is visible in the ranking.} Niazimicin holds
the only full identity score among compounds confirmed in a dichloromethane fraction
($D_1=2$) and scores \emph{zero} on radiation biology, so it sits at tier~C. Moringin, the
signature \textit{Moringa} isothiocyanate, reaches only tier~B. Four of the five tier-A
compounds hold $D_1=1$, a spectral-library or dereplication match, and none of the five has
been found in a dichloromethane extract.

\section{What the published biological evidence shows}\label{sec:biology}

The curated evidence establishes the following.
\begin{itemize}
  \item \textbf{Moringin, the isothiocyanate most characteristic of \textit{Moringa}, has no
    ionising-radiation study.} Exhaustive searching found no primary study combining
    moringin or glucomoringin with ionising radiation. Its evidence is confined to work
    without radiation: pro-apoptotic effects involving p53, CDKN1A, BAX and caspases in
    neuroblastoma cells \parencite{cirmi2019}, G2/M arrest in prostate cancer cells
    \parencite{abdkarim2023}, NRF2 activation in a motor-neuron model
    \parencite{gugliandolo2024}, and NF-$\kappa$B inhibition with a myeloma xenograft
    \parencite{brunelli2010}. Because chemical-class analogy earns no points, moringin gains
    nothing from the benzyl isothiocyanate radiation literature, and that analogy is
    treated as an untested hypothesis.
  \item \textbf{The best-evidenced radiation candidates are not \textit{Moringa}-specific.}
    $\beta$-Sitosterol has normal-tissue radioprotection evidence from two independent
    groups \parencite{li2007,elbakary2026}, as does rutin \parencite{patil2017,thabet2018},
    although one research group contributes to both compounds' evidence;
    $\alpha$-tocopherol has a rat crypt-microcolony clonogenic result
    \parencite{felemovicius1995} and a randomised controlled trial in radiation mucositis
    \parencite{ferreira2004}. All three are documented in \textit{Moringa} seed or leaf, but
    \textbf{none has been found in a dichloromethane extract}, which is why the extraction
    evidence is shown beside the score rather than within it.
  \item \textbf{One candidate carries an explicit liability.} Chlorogenic acid has
    normal-tissue protection evidence in bone marrow \parencite{abraham1993} and testis
    \parencite{abedpour2022}, but \textcite{yin2022} report that it \emph{protects
    hepatocellular carcinoma cells} from radiation through NRF2 and conclude that it could
    hinder radiotherapy. Under the prespecified rule this caps its tier at C. Astragalin is
    capped for a related reason: its only DNA-damage-relevant primary signal is
    \emph{pro-genotoxic} \parencite{saito2012}, and the one astragalin \enquote{radiation}
    paper concerns UVB, which is not ionising radiation.
  \item \textbf{An informative negative was retained.} Membrane oleate enrichment did
    \textbf{not} alter the clonogenic radiosensitivity of human lymphoid cells
    \parencite{george1983}. This direct clonogenic test argues against prioritising oleic
    acid despite its abundance in seed oil.
\end{itemize}

\begin{figure}[p]
  \centering
  \includegraphics[width=\linewidth,height=0.62\textheight,keepaspectratio]{fig_08_compound_priority}
  \caption[Evidence-based prioritisation of \textit{Moringa} seed and leaf
    constituents]{Evidence-based prioritisation of documented \textit{Moringa oleifera} seed
    and leaf constituents for future experimental work. (a)~Points awarded in each of the
    four prespecified dimensions (numbers are points; shading is the fraction of the
    dimension's maximum), beside the extraction-specific evidence: identified in a
    dichloromethane extract or fraction of seed or leaf (dark), in seed or leaf with another
    solvent (light), or only indirectly (white). (b)~Total priority score out of eight,
    coloured by tier. The extraction column is kept outside the score because the original
    study used dichloromethane extracts whose composition was never measured. The scoring
    rules were fixed before any score was computed; every tier, including A, is
    hypothesis-generating, and no tier asserts radioprotection, radiosensitisation, potency
    or any effect of the original extracts. Two systematic names are shortened in the
    figure (see the note to Table~\ref{tab:ranking}).}
  \label{fig:priority}
\end{figure}

\paragraph{Reading Figure~\ref{fig:priority}.} Down the extraction column and across the
four dimensions, the \textit{Moringa}-specific compounds with the darkest identity cell ---
niazimicin, glucomoringin, niazirin, niaziminin, the acetylated isothiocyanate and the
carbamate glycoside, all with a structurally confirmed presence in seed or leaf --- are
precisely those whose radiation-biology cell is empty or nearly so. The one compound with
full scores for both identity and radiation biology is $\alpha$-tocopherol, a ubiquitous
plant metabolite that has not been found in a dichloromethane extract; the other four
tier-A compounds carry only one identity point. Panel~(b) shows the effect of the liability
rule: chlorogenic acid's total of 7 equals that of four tier-A compounds, yet its bar is
coloured as tier~C. A reader comparing bar lengths alone would place it among the leaders;
the prespecified cap prevents that, and the figure shows the discrepancy rather than hiding
it.

\begin{figure}[p]
  \centering
  \includegraphics[width=\linewidth,height=0.62\textheight,keepaspectratio]{fig_09_compound_evidence_map}
  \caption[What the published evidence links each shortlisted compound to]{What the
    published evidence links each shortlisted compound to. (a)~Radiation-response processes
    for which a primary study of the purified compound exists: symbol size is the number of
    independent research groups, and colour distinguishes studies using ionising radiation
    from studies of the same process without radiation. Compounds are in priority order. An
    empty row means that no qualifying primary study was located, which is a gap in the
    literature, not evidence of no effect. (b)~Number of genes measured in a published
    experiment on the compound that are also members of a Hallmark gene set ranked by the
    exploratory explainability analysis (Table~\ref{tab:c4}); blank cells mean no overlap was
    recorded. An overlap marks a point where a compound and a model-highlighted gene set
    touch the same genes. It is a hypothesis to be tested, not evidence that the compound
    modifies the gene set, the pathway or clonogenic SF2, and the gene-set ranking itself is
    capped at tier~C. Reported directions of effect are given in Table~\ref{tab:c2} and are
    never merged.}
  \label{fig:evidence-map}
\end{figure}

\paragraph{Reading Figure~\ref{fig:evidence-map}.} Attributed compound by compound, the
biological evidence is thin. Only a minority of cells in panel~(a) are filled, and several
shortlisted compounds have no ionising-radiation evidence on their row: moringin, the most
\textit{Moringa}-characteristic candidate, appears only in the non-radiation colour, and
niazimicin has no scored row at all because its one biological study dosed an extract
fraction rather than the purified compound. Panel~(b) is deliberately sparse. Only 25
compound-by-gene-set overlaps were recorded (Table~\ref{tab:c4}), 23 for shortlisted
compounds and two for benzyl isothiocyanate, which is not assessable and is not shown. They
are small: 14 involve a single gene and 9 involve two, and the largest are moringin with the
p53 pathway (three genes: BAX, CDKN1A, TP53) and quercetin with the p53 pathway (five genes:
BAX, CDKN1A, DDIT3, TGFB1, TP53). These overlaps are the thinnest kind of link in the
project, and they mark where a hypothesis could be formed, not where an effect has been
shown.

\section{What a tier means, and what it does not}\label{sec:tier-meaning}

A tier records \textbf{the state of the published evidence about a compound}, assembled under
rules fixed before any score was computed (Table~\ref{tab:tier-meaning}). It is not a measure
of radioprotection. Tier~A means \enquote{best documented}, not \enquote{most likely to
work}. Four of the five tier-A compounds are ubiquitous plant metabolites whose radiation
evidence comes entirely from outside \textit{Moringa} research, and the direction of that
evidence differs between them: quercetin and kaempferol have \emph{radiosensitisation}
evidence in tumour cells as well as normal-tissue protection, whereas $\beta$-sitosterol and
rutin have normal-tissue protection only. Radioprotection and radiosensitisation are
opposite effects and are never merged (Table~\ref{tab:c2}).

\begin{table}[tbp]
  \caption{What a compound tier does and does not say.}\label{tab:tier-meaning}
  \small
  \begin{tabular}{@{}L{0.47\linewidth}L{0.47\linewidth}@{}}
    \toprule
    \thead{A tier does say} & \thead{A tier does not say} \\
    \midrule
    Primary studies of the purified compound exist, with the stated independence
      & That the compound protects or sensitises anything \\
    The compound is documented in \textit{Moringa} seed or leaf at the stated
      identification strength & That it was present in the extracts of the original study \\
    Published work connects it to a radiation-response process
      & That it modifies clonogenic survival, the endpoint modelled here \\
    A follow-up experiment can be specified from the record
      & That such an experiment would succeed \\
    \bottomrule
  \end{tabular}
\end{table}

\section{Realistic first-stage experimental priorities}\label{sec:first-stage}

The ranking answers the question \enquote{what is best documented?}. A first experiment has to
answer a narrower question, so the two orderings differ --- without any score being changed.
Three questions can be separated.

\paragraph{Is any plausibly active compound present in the extracts that were tested?}
This is the project's largest gap and the cheapest to close. The priority candidates are
\textbf{niazimicin} (tier~C, the only compound with a confirmed structural identification in
both the seed and the leaf dichloromethane fractions) and \textbf{moringin} (tier~B, the
signature \textit{Moringa} isothiocyanate). Neither has any ionising-radiation study, so
neither can be offered as a protection candidate; they are priorities for
\textbf{analytical characterisation of the original extracts}, followed by a first radiation
screen if they are found.

\paragraph{Does a compound with existing radiation evidence act on the endpoint modelled
here?} \textbf{$\alpha$-Tocopherol} is the strongest first-stage candidate: it is the only
tier-A compound with a full identity score, and the only compound in the evidence base whose
normal-tissue protection evidence includes a \emph{clonogenic} endpoint (rat intestinal crypt
survival), alongside a randomised trial in radiation mucositis. \textbf{$\beta$-Sitosterol} follows,
with normal-tissue protection reported by two independent groups. Both must be tested in
matched normal and tumour cells, because neither has a tumour arm.

\paragraph{Which candidates should not lead?} \textbf{Chlorogenic acid} should not be a
first-stage protection candidate despite its score: the tumour-protection finding makes it a
cautionary compound in a radiotherapy context. \textbf{Astragalin} is capped for a related
reason, its only DNA-damage-relevant evidence being pro-genotoxic. \textbf{Oleic acid} has a
direct clonogenic test showing \emph{no} effect on radiosensitivity, so it is informative but
not a candidate. Quercetin, kaempferol and rutin are scientifically interesting but
second-stage: their \textit{Moringa} identity rests on tentative dereplication, and the
aglycone--glycoside pharmacokinetic question would have to be settled first.

\begin{keybox}[title={A defensible first-stage programme.}]
Characterise the original extracts (niazimicin, moringin), and test $\alpha$-tocopherol and
$\beta$-sitosterol on clonogenic survival in matched normal and tumour cells over a dose
range. This uses the tiers as summaries of evidence, not as predictions, and leaves every
score exactly as computed.
\end{keybox}

The scientific shape of this module's result is therefore a \textbf{mismatch}, and it is its
most useful output: the compound with the best claim to being present in the tested extracts
has no radiation biology, and the compounds with radiation biology have no confirmed presence
in those extracts. Neither gap can be closed by further literature work; both can be closed by
the experiments proposed in Chapter~\ref{ch:recommendations}.
